\begin{derivation}[从\eqnrefs{eq:virtual-path-neumann-dp85}到\eqnrefs{eq:virtual-path-remainder-bound-dp85}]
在 $Q$ 子空间内写
\begin{equation*}
 E-QHQ
 =E-QH_0Q-QVQ
 =[1-G_Q^{(0)}QVQ]\,[G_Q^{(0)}]^{-1}.
\end{equation*}
因此
\begin{equation*}
 (E-QHQ)^{-1}
 =[1-G_Q^{(0)}QVQ]^{-1}G_Q^{(0)}.
\end{equation*}
若 $\epsilon_Q=\|G_Q^{(0)}QVQ\|<1$，Neumann 级数在算符范数下收敛：
\begin{align*}
 [1-G_Q^{(0)}QVQ]^{-1}
 &=\sum_{r=0}^{\infty}(G_Q^{(0)}QVQ)^r,\\
 (E-QHQ)^{-1}
 &=\sum_{r=0}^{\infty}(G_Q^{(0)}QVQ)^rG_Q^{(0)},
\end{align*}
即正文\eqnrefs{eq:virtual-path-neumann-dp85}。

把该结果代入
\begin{equation*}
 H_{\rm eff}(E)=PHP+PHQ(E-QHQ)^{-1}QHP
\end{equation*}
并在 $P$ 空间两态之间取矩阵元：
\begin{align*}
 \langle b|H_{\rm eff}|a\rangle
 ={}&\langle b|PHP|a\rangle\\
 &+\sum_{r=0}^{\infty}
 \langle b|PHQ
 (G_Q^{(0)}QVQ)^rG_Q^{(0)}QHP|a\rangle,
\end{align*}
得到正文\eqnrefs{eq:virtual-path-general-series-dp85}。

若只保留 $r=0,\ldots,R$，预解式余项为
\begin{equation*}
 \delta G_Q^{(>R)}
 =\sum_{r=R+1}^{\infty}(G_Q^{(0)}QVQ)^rG_Q^{(0)}.
\end{equation*}
利用算符范数的次乘性与几何级数，
\begin{align*}
 \|\delta G_Q^{(>R)}\|
 &\le
 \sum_{r=R+1}^{\infty}\epsilon_Q^r\,\|G_Q^{(0)}\|\\
 &=\|G_Q^{(0)}\|
 \frac{\epsilon_Q^{R+1}}{1-\epsilon_Q}.
\end{align*}
再在左右乘上 $PHQ$ 与 $QHP$，得到正文\eqnrefs{eq:virtual-path-remainder-bound-dp85}。
\end{derivation}

\begin{derivation}[从\eqnrefs{eq:virtual-path-general-series-dp85}到\eqnrefs{eq:nth-bragg-ac-stark-dp83}]
取正文已经定义的最近邻虚链。自由部分在格点基底中对角：
\begin{equation*}
 H_0|j\rangle=E_j|j\rangle,
\end{equation*}
相互作用只连接最近邻：
\begin{equation*}
 V
 =\sum_j\frac{\hbar\Omega_j}{2}
 \left(|j+1\rangle\langle j|+|j\rangle\langle j+1|\right).
\end{equation*}
目标子空间与补空间为
\begin{align*}
 P&=|0\rangle\langle0|+|N\rangle\langle N|,\\
 Q&=\sum_{j=1}^{N-1}|j\rangle\langle j|+Q_{\rm ext},
\end{align*}
其中 $Q_{\rm ext}$ 表示目标区间之外的其他格点；计算最短 $0\to N$ 路径时它们不参与最低阶非对角矩阵元。取展开能量 $E_*$，定义
\begin{equation*}
 \Delta_j=\frac{E_*-E_j}{\hbar},
\end{equation*}
则未扰动 $Q$ 预解式在内部格点上严格为
\begin{equation*}
 G_Q^{(0)}(E_*)|j\rangle
 =\frac1{E_*-E_j}|j\rangle
 =\frac1{\hbar\Delta_j}|j\rangle.
\end{equation*}

先求非对角元 $\langle N|H_{\rm eff}|0\rangle$。由于 $PVP$ 对 $N>1$ 为零，最低非零项必须含 $N$ 个最近邻顶角和 $N-1$ 个中间传播子。正文\eqnrefs{eq:virtual-path-general-series-dp85} 中对应的项是
\begin{align*}
 \Sigma_{N0}^{(N)}(E_*)
 ={}&\langle N|VQG_Q^{(0)}QVQG_Q^{(0)}\cdots
 QVQG_Q^{(0)}QV|0\rangle,
\end{align*}
其中共有 $N-2$ 个内部 $QVQ$，加上两端的 $PVQ$ 与 $QVP$，总顶角数为 $N$。在最近邻假设下，每插入一次恒等分解
\begin{equation*}
 Q=\sum_{j=1}^{N-1}|j\rangle\langle j|+Q_{\rm ext}
\end{equation*}
以后，所有不满足相邻格点选择律的项都为零，因此最低阶只剩唯一单调路径
\begin{equation*}
 0\to1\to2\to\cdots\to N-1\to N.
\end{equation*}
逐个写出矩阵元：
\begin{align*}
 \Sigma_{N0}^{(N)}(E_*)
 ={}&\langle N|V|N-1\rangle
 \frac1{E_*-E_{N-1}}
 \langle N-1|V|N-2\rangle\\
 &\times\frac1{E_*-E_{N-2}}
 \cdots
 \frac1{E_*-E_1}
 \langle1|V|0\rangle.
\end{align*}
利用
\begin{equation*}
 \langle j+1|V|j\rangle=\frac{\hbar\Omega_j}{2},
 \qquad
 E_*-E_j=\hbar\Delta_j,
\end{equation*}
得到
\begin{align*}
 \Sigma_{N0}^{(N)}(E_*)
 &=\frac{(\hbar\Omega_{N-1}/2)\cdots(\hbar\Omega_0/2)}
 {(\hbar\Delta_{N-1})\cdots(\hbar\Delta_1)}\\
 &=\frac{\hbar}{2}
 \frac{\Omega_{N-1}\Omega_{N-2}\cdots\Omega_0}
 {2^{N-1}\Delta_{N-1}\Delta_{N-2}\cdots\Delta_1}.
\end{align*}
按正文采用的二态非对角元约定
\begin{equation*}
 \langle N|H_{\rm eff}|0\rangle
 \equiv\frac{\hbar\Omega_{\rm eff}^{(N)}}2,
\end{equation*}
立即得到正文\eqnrefs{eq:nth-bragg-effective-rabi-dp83} 的首项。

为什么下一项比首项多两个顶角也可以从路径结构直接看出。最近邻格上，从 $0$ 到 $N$ 的路径净向右位移必须为 $N$。若额外增加一步向左，则为了最终仍到达 $N$，必须再增加一步向右。因此相对于最短 $N$ 顶角路径，任何更长路径至少含 $N+2$ 个顶角。若所有耦合满足 $|\Omega_j|\le\Omega_*$，所有内部能量分母满足 $|\Delta_j|\ge\Delta_*$，则下一阶非对角矩阵元的尺度至多为
\begin{equation*}
 \frac{|\delta\Sigma_{N0}|}{\hbar}
 =O\!\left(
 \frac{\Omega_*^{N+2}}{\Delta_*^{N+1}}
 \right),
\end{equation*}
对应正文\eqnrefs{eq:nth-bragg-effective-rabi-dp83} 的修正阶数；如果存在更远程连边，这个“多两个顶角”的计数不再成立，应改用一般算符范数余项界~\eqnrefs{eq:virtual-path-remainder-bound-dp85}。

再求端点的最低阶对角自能。对任一保留态 $|a\rangle\in P$，Feshbach 自能的二阶项为
\begin{equation*}
 \delta E_a^{(2)}
 =\langle a|VQG_Q^{(0)}Q V|a\rangle.
\end{equation*}
在 $Q$ 子空间插入完整本征态分解，不能预先丢掉链外邻态：
\begin{align*}
 \delta E_a^{(2)}
 &=\sum_{m\in Q}
 \langle a|V|m\rangle
 \frac1{E_*-E_m}
 \langle m|V|a\rangle\\
 &=\sum_{m\in Q}
 \frac{|\langle m|V|a\rangle|^2}{E_*-E_m}.
\end{align*}
只对非零矩阵元保留求和项，并用正文记号
\begin{equation*}
 \langle m|V|a\rangle=\frac{\hbar\Omega_{ma}}2,
 \qquad
 \Delta_m^{(a)}=\frac{E_*-E_m}{\hbar},
\end{equation*}
便有
\begin{align*}
 \delta E_a^{(2)}
 &=\sum_{m\in Q:\,\langle m|V|a\rangle\ne0}
 \frac{\hbar^2|\Omega_{ma}|^2/4}{\hbar\Delta_m^{(a)}}\\
 &=\hbar
 \sum_{m\in Q:\,\langle m|V|a\rangle\ne0}
 \frac{|\Omega_{ma}|^2}{4\Delta_m^{(a)}}.
\end{align*}
除以 $\hbar$ 即得到正文\eqnrefs{eq:nth-bragg-ac-stark-dp83}。若链在两端真正截断，则 $a=0$ 的求和只剩 $m=1$，$a=N$ 的求和只剩 $m=N-1$，从而分别约化为
\begin{equation*}
 \delta_0^{\rm AC}=\frac{|\Omega_0|^2}{4\Delta_1},
 \qquad
 \delta_N^{\rm AC}=\frac{|\Omega_{N-1}|^2}{4\Delta_{N-1}}.
\end{equation*}
若目标态位于无限最近邻格中，则还存在链外邻态，必须把相应 $m$ 项保留在同一求和式中。

最后估计中间态泄漏。对最邻近中间态，把除 $0\leftrightarrow1$ 直接耦合以外的所有项记作 $r_1$，定态方程严格写成
\begin{equation*}
 (E_*-E_1)c_1
 +\frac{\hbar\Omega_0}{2}c_0+r_1=0.
\end{equation*}
由 $\Delta_1=(E_*-E_1)/\hbar$ 得到恒等式
\begin{equation*}
 c_1
 =-\frac{\Omega_0}{2\Delta_1}c_0
 -\frac{r_1}{\hbar\Delta_1}.
\end{equation*}
Neumann 虚路径展开给 $r_1/(\hbar\Delta_1)=O[(\Omega_*/\Delta_*)^2]$，所以首项振幅为 $O(\Omega_*/\Delta_*)$，其概率从 $O[(\Omega_*/\Delta_*)^2]$ 开始；其他中间态沿连续虚路径继续增加额外的 $\Omega/\Delta$ 因子。
\end{derivation}
